% ─────────────────────────────────────────────────────────────────────────────
This chapter reviews the literature across five areas that underpin BookForMe:
agentic AI systems; task-oriented dialogue and state tracking; bilingual and
code-switched NLU; model selection and training for low-resource settings; and
distributed booking systems with concurrency guarantees. Each section motivates
the specific technical decisions made in the project and identifies the gaps in
existing work that this project addresses.

% ─────────────────────────────────────────────────────────────────────────────
\section{From Chatbots to Autonomous Agents}
\label{sec:lit-agents}

\subsection{Defining the Agentic Architecture}

The term ``agent'' has been used broadly in software engineering and robotics,
but its application to large language models (LLMs) represents a specific
architectural shift. Wang et al.~\cite{wang2024} propose a unified framework in
which an LLM-based agent comprises four modules: \textit{Profiling} (role and
persona definition), \textit{Memory} (short-term context and long-term
persistent state), \textit{Planning} (decomposing goals into subgoals and
choosing actions), and \textit{Action} (calling external tools and APIs to
effect changes in the environment). The LLM functions as a coordinator that
orchestrates a Perception $\rightarrow$ Reasoning $\rightarrow$ Action cycle,
fundamentally distinguishing it from a passive text predictor.

Xi et al.~\cite{xi2023} further distinguish agents from traditional
conversational systems by their capacity to \textit{initiate and manage
workflows} rather than merely respond to prompts. A conventional chatbot can
answer ``Is the 5~PM slot available?''; our agent must answer ``Book the 5~PM
slot, lock it for ten minutes, verify the advance payment, and notify the
vendor.'' The latter requires stateful read-write operations across multiple
external systems, which is precisely the capability gap this work fills.

\subsection{Cognitive Architectures: ReAct, PRACT, and Reflexion}

Yao et al.~\cite{yao2023} introduced the \textbf{ReAct} (Reasoning + Acting)
paradigm, which interleaves chain-of-thought reasoning with external actions,
enabling explicit planning traces such as:\\
\textit{Thought}: User requests 5~PM padel.
$\rightarrow$ \textit{Action}: \texttt{get\_availability("padel","2026-03-15","17:00")}.
$\rightarrow$ \textit{Observation}: Slot available.
$\rightarrow$ \textit{Action}: \texttt{create\_booking(user, slot)}.

Liu et al.~\cite{liu2024} extend this with the \textbf{PRACT} agent, adding a
self-reflection step before executing irreversible actions. Before calling
\texttt{create\_booking()}, the agent verifies that the action is consistent
with the user's latest stated intent, which is a critical safeguard where a
confirmed booking constitutes a financial commitment.

Shinn et al.~\cite{shinn2023} demonstrate in \textbf{Reflexion} that verbal
reinforcement and persistent state across iterations significantly improve agent
reliability on multi-step tasks. The agent receives textual feedback from the
environment (e.g., ``payment amount incorrect'') and updates its policy
accordingly without gradient updates. This pattern directly informs our OCR
rejection loop, where the agent re-prompts the user upon detecting an incorrect
payment amount.

\subsection{Orchestration Frameworks: LangChain and LangGraph}

LangChain~\cite{mavroudis2024} provides a modular abstraction layer that
decouples application logic from specific LLM providers, enabling composition
of prompt templates, vector stores, and output parsers into unified pipelines.
However, standard LangChain implementations rely on Directed Acyclic Graphs
(DAGs), which cannot represent the iterative, self-correcting loops required
for asynchronous WhatsApp conversations where users may reply late, change
their mind, or submit incorrect payment proofs.

\textbf{LangGraph}~\cite{langgraph2024} addresses this by replacing DAGs with
\textit{Cyclic State Graphs}: the agent's behaviour is modelled as a state
machine where nodes represent actions and edges encode conditional logic. The
framework provides loop control, state persistence across invocations, and
fault recovery, all of which are essential for BookForMe's multi-turn WhatsApp
booking conversations. Shang et al.~\cite{shang2024} confirm in AgentSquare
that modular architectures separating Planning, Reasoning, and Memory are
substantially more robust for complex tasks than monolithic prompt designs.

\subsection{Multi-Agent Decomposition}

H{\"a}ndler~\cite{handler2023} proposes a multi-dimensional taxonomy of
autonomous LLM-powered multi-agent architectures, showing that complex
workflows benefit from assigning distinct roles to specialised sub-agents
(intent parsing, verification, database updates, external communications)
coordinated by a central orchestrator. Wu et al.~\cite{wu2023autogen} in
AutoGen further argue that effective problem-solving requires cyclic loops: the
ability to retry failed actions, request missing information, and self-correct
rather than proceeding through purely linear pipelines. BookForMe adopts this
multi-node, cyclic design in its LangGraph implementation.

% ─────────────────────────────────────────────────────────────────────────────
\section{Task-Oriented Dialogue and State Tracking}
\label{sec:lit-tod}

Task-Oriented Dialogue (TOD) systems have evolved from modular pipelines
(NLU $\to$ Dialogue State Tracking $\to$ Policy $\to$ NLG) to end-to-end
architectures. \textbf{SimpleTOD}~\cite{hosseini2020} treats the entire
dialogue as unified causal language modelling, jointly predicting belief
states, system actions, and responses in a single autoregressive model. While
effective on standard benchmarks (MultiWOZ), it requires large annotated
corpora, which is a significant barrier for our domain where no pre-built
dataset exists.

Shin et al.~\cite{shin2022} address this data scarcity problem with
\textbf{DS2} (Dialogue Summaries as Dialogue States), converting conversation
history into templated natural-language summaries that map directly to
structured slot representations. This summarisation-based approach is
particularly useful for short, noisy WhatsApp messages where maintaining a
full belief-state vector is impractical.

\subsection{Why Separate NLU and NLG}

End-to-end TOD systems that jointly handle understanding and generation perform
well when large in-domain corpora exist. Peng et al.~\cite{peng2021} in
\textbf{Soloist} demonstrate that a unified generative model can handle full
dialogue turns at scale through transfer learning and machine teaching, but
this approach still requires thousands of task-specific annotated dialogues
and introduces non-determinism in every response, including slot fills and
booking decisions. For a production booking system where incorrect slot data
leads to real financial transactions, non-determinism in the understanding
stage is unacceptable.

This motivates BookForMe's two-model split: a deterministic, fine-tuned
classifier (DistilBERT) for intent routing, and a large generative model
(Qwen 3 32B via Groq) only for surface-level natural language generation of
responses. The classifier's output is typed and bounded; the generator never
touches database state directly. This is consistent with the Neuro-Symbolic
architecture recommended by H{\"a}ndler~\cite{handler2023} and Wu et
al.~\cite{wu2023autogen}.

% ─────────────────────────────────────────────────────────────────────────────
\section{Bilingual and Code-Switched NLU}
\label{sec:lit-bilingual}

\subsection{Roman Urdu Challenges}

Processing Roman Urdu introduces challenges largely absent from standard
multilingual NLP. Bhat et al.~\cite{bhat2023} classify Roman Urdu as
exhibiting severe \textit{orthographic variation}: a single word can be
rendered in multiple phonetically plausible romanisations by different writers
(e.g., \textit{waqt}, \textit{vakt}, \textit{wkht}, \textit{waqat}, and
\textit{tym} all mean ``time''). Standard subword tokenisers (BPE, WordPiece)
trained on formal corpora frequently fail to map these variants to a shared
semantic representation. A normalisation step or a model pre-exposed to
informal text is therefore necessary before intent classification.

\subsection{Cross-Lingual Transfer and Adapters}

Yu et al.~\cite{yu2023} address cross-lingual dialogue state tracking via
contrastive learning, aligning slot embeddings across languages so that
semantically equivalent slots in different languages map to nearby vector
representations. Wu et al.~\cite{wu2023adapter} show that parameter-efficient
adapter methods can transfer pretrained models to code-switched contexts
without full retraining, significantly reducing data and compute requirements.

The Qwen2.5 technical report~\cite{qwen2025} demonstrates strong multilingual
instruction-following capability, validating its use for generating bilingual
booking responses in English and Roman Urdu where tone, code-switching, and
natural phrasing are all important.

\subsection{Pretrained Models for Classification}

Devlin et al.~\cite{devlin2019} introduced \textbf{BERT}, establishing
bidirectional masked language modelling as the dominant pretraining paradigm.
The multilingual variant (mBERT) was pretrained on 104 languages including
Urdu, providing limited but useful cross-lingual transfer. Conneau et
al.~\cite{conneau2020} demonstrate that \textbf{XLM-RoBERTa}, trained on a
substantially larger multilingual corpus with improved data filtering,
outperforms mBERT on cross-lingual classification tasks.

Sanh et al.~\cite{sanh2019} show that \textbf{DistilBERT}, a distilled
66M-parameter version of BERT, retains 97\% of BERT's NLP benchmark
performance while being 40\% smaller and 60\% faster at inference. Knowledge
distillation~\cite{hinton2015} achieves this by training a compact student
model to replicate the soft output distribution of a larger teacher, preserving
most of the representational quality at a fraction of the compute cost. For a
latency-sensitive production system like BookForMe, where intent classification
must complete in under 50~ms per message, this speed-accuracy tradeoff is the
primary motivation for choosing DistilBERT over larger alternatives.

\subsection{Data Augmentation for Low-Resource Settings}

Given the absence of any pre-built bilingual booking dataset, data augmentation
is essential. Hou et al.~\cite{hou2020} show that back-translation and
contextual augmentation improve slot-filling F1 by 6--17\% and intent accuracy
by up to 12\% in booking domains with fewer than 500 annotated utterances. An
et al.~\cite{an2022} demonstrate that controllable T5-based generation yields
8--18\% absolute gains in intent classification under 10-shot conditions for
restaurant and hotel reservation tasks. He et al.~\cite{he2024} in
\textbf{SynthDST} show that fully synthetic LLM-generated dialogues can
surpass real-data baselines, achieving 4--12\% higher Joint Goal Accuracy on
MultiWOZ when human annotations are extremely limited. BookForMe's augmentation
strategy follows these findings directly.

Focal loss~\cite{lin2017}, originally introduced for object detection, is
effective for training on class-imbalanced datasets: it down-weights easy,
well-classified examples and focuses training signal on the harder class
boundaries. Given the imbalance in our 5-class corpus (39.9\% inquiry vs.
5.3\% greeting), focal loss is more appropriate than standard cross-entropy.

% ─────────────────────────────────────────────────────────────────────────────
\section{Retrieval, Planning, and Tool Use}
\label{sec:lit-tools}

Schick et al.~\cite{schick2023} demonstrated in \textbf{Toolformer} that LLMs
can learn to invoke external APIs through self-supervised training, deciding
autonomously when and how to call tools. Patil et al.~\cite{patil2023}
(\textbf{Gorilla}) and Qin et al.~\cite{qin2023} (\textbf{ToolLLM}) extend
this to controlled invocation of structured APIs, providing the theoretical
basis for BookForMe's \texttt{get\_availability()},
\texttt{create\_booking()}, and \texttt{get\_services()} tool set.

Lewis et al.~\cite{lewis2020} introduced \textbf{Retrieval-Augmented
Generation (RAG)}, conditioning generation on dynamically retrieved documents.
In BookForMe, the agent retrieves up-to-date vendor availability and pricing
from Firestore rather than relying on memorised parametric knowledge, which is
critical for a real-time booking system. Chen et al.~\cite{chen2025} present
adaptive retrieval strategies where the agent decides when retrieval is
necessary, minimising unnecessary database queries and reducing latency.

% ─────────────────────────────────────────────────────────────────────────────
\section{Payment Verification via Computer Vision}
\label{sec:lit-ocr}

A distinguishing feature of BookForMe is its automated payment screenshot
verification via OCR, eliminating the need for manual receipt checking by
vendors. Memon et al.~\cite{memon2020} provide a comprehensive review of OCR
systems, noting that modern deep learning-based OCR achieves high accuracy on
printed and near-printed text such as bank transfer screenshots, where fonts
are consistent and the background is predictable. For edge cases such as
blurry images or unusual receipt formats, their analysis recommends requesting
a cleaner image rather than forcing a low-confidence extraction, which aligns
exactly with our fallback strategy.

The Gemini multimodal model~\cite{team2023gemini} extends this capability to
structured prompting: given a screenshot and a question (``what amount was
transferred?''), the model returns a specific extracted value rather than a
general description. This approach is more reliable for our use case than
training a dedicated OCR model, as it leverages a large pretrained vision
backbone without requiring any labelled payment screenshot data.

% ─────────────────────────────────────────────────────────────────────────────
\section{Distributed Booking and Concurrency Safety}
\label{sec:lit-concurrency}

Mong'are~\cite{mongare2024} provides a thorough analysis of idempotency in
APIs and distributed systems, specifically addressing the double-booking
problem that arises when concurrent requests modify shared mutable state. The
recommended approach is idempotent APIs combined with Optimistic Concurrency
Control (OCC) or distributed locking. BookForMe's implementation directly
follows this recommendation, using Firestore atomic transactions with a
\texttt{hold\_expires\_at} timestamp to implement OCC.

Singh et al.~\cite{singh2000} demonstrate the effectiveness of structured
state management for dialogue policy in the NJFun system, validating the
importance of separating dialogue state from the booking transaction state,
which is exactly the separation BookForMe maintains between the LangGraph
\texttt{AgentState} and the Firestore slot document.

% ─────────────────────────────────────────────────────────────────────────────
\section{Social Ranking: Elo Rating Systems}
\label{sec:lit-elo}

The BookForMe social layer includes a competitive matchmaking system using Elo
ratings. The Elo system~\cite{elo1978}, originally developed for chess, assigns
each player a numerical rating updated after each match based on the expected
versus actual outcome, with the magnitude of the update weighted by the
opponent's relative strength. This produces a self-calibrating ranking that
accounts for the quality of opponents rather than simply counting wins. It is
widely adopted in competitive online games (e.g., League of Legends, Chess.com)
and is well understood by sports players, making it an appropriate choice for
the casual and competitive match lobbies in BookForMe's social hub.

% ─────────────────────────────────────────────────────────────────────────────
\section{Novelty and Gap Analysis}
\label{sec:lit-gap}

Table~\ref{tab:lit-gap} summarises the specific gaps that BookForMe addresses
relative to the existing literature and systems.

\begin{table}[H]
\centering
\caption{Gap analysis: BookForMe versus existing dialogue and booking paradigms.}
\label{tab:lit-gap}
\small
\begin{tabularx}{\textwidth}{>{\raggedright\arraybackslash}p{0.20\textwidth}YC{0.10\textwidth}C{0.12\textwidth}C{0.12\textwidth}C{0.12\textwidth}}
\toprule
\textbf{System} & \textbf{Focus} & \textbf{Tool Use} &
\textbf{Roman Urdu} & \textbf{Async Mem.} & \textbf{Txn.\ Safety} \\
\midrule
GPT / General LLM    & General chat        & \cmark & \cmark (partial) & \xmark & \xmark \\
SimpleTOD            & Research datasets   & \xmark & \xmark           & \xmark & \xmark \\
ReAct Agents         & Web automation      & \cmark & \xmark           & \cmark (limited) & \xmark \\
Domain-specific apps & Single category     & \xmark & \xmark           & \xmark & \cmark (partial) \\
\textbf{BookForMe}   & \textbf{WA Booking} & \cmark & \cmark           & \cmark & \cmark \\
\bottomrule
\end{tabularx}
\end{table}

The primary research gaps addressed are:
\begin{itemize}
  \item \textbf{Roman Urdu code-switching:} Existing DST and NLU research does
        not target informal WhatsApp-style Roman Urdu, presenting gaps in
        language modelling and orthographic normalisation for this dialect.
  \item \textbf{Transactional booking guarantees:} While agent tool-use is well
        studied, very few works address concurrency, idempotency, and
        transactional safety in real-world booking systems.
  \item \textbf{Asynchronous conversation flow:} Most TOD frameworks assume
        synchronous interaction. WhatsApp requires persistent memory and session
        resumption across arbitrary time delays.
  \item \textbf{Evaluation for informal-economy SMEs:} Agent evaluation
        literature does not address booking systems in informal economies where
        reliability, low latency, and no-infrastructure vendor onboarding are
        essential constraints.
\end{itemize}

BookForMe synthesises the reviewed literature into a practical LLM-driven
booking agent that is novel precisely because it operates at the intersection
of all four gaps listed above.

%%% Local Variables:
%%% mode: latex
%%% TeX-master: "../report"
%%% End:
